\chapter[Conclusiones]{Conclusiones} \label{cp:conclusiones}

A lo largo de este trabajo se abordó el diseño, la simulación, la fabricación y la caracterización experimental de un amplificador de potencia Clase AB en banda S para aplicaciones de enlace descendente en CubeSats, junto con un acoplador híbrido en cuadratura concebido como primera etapa hacia una futura arquitectura Doherty. El desarrollo completo, desde el marco teórico hasta la medición de los prototipos fabricados, permitió recorrer la totalidad del proceso de diseño de un sistema de RF para aplicaciones espaciales, contemplando tanto las restricciones propias de cada circuito como las impuestas por la plataforma CubeSat y por el entorno de órbita baja terrestre.

El acoplador híbrido en cuadratura, implementado mediante una topología de línea ramificada plegada para reducir sus dimensiones, alcanzó en su segunda versión una frecuencia central de 2,019~GHz y un desfasaje de 92,57° entre los puertos de salida en la frecuencia de interés, valores considerablemente más próximos al comportamiento ideal que los obtenidos con la primera versión fabricada. Esta mejora entre versiones validó de manera directa la hipótesis formulada durante el proceso de diseño, según la cual el corrimiento de frecuencia observado se debía a la falta de modelado de la transición hacia el conector SMA, y confirmó la utilidad de la verificación electromagnética de onda completa como herramienta de diagnóstico frente a las limitaciones del modelo esquemático de líneas de transmisión ideales. La comparación entre los resultados medidos y los obtenidos mediante simulación en ADS y en CST mostró, además, una concordancia satisfactoria entre ambas plataformas de simulación y el prototipo fabricado, con una mayor precisión de CST atribuible a su capacidad de capturar los efectos parásitos de la geometría real del layout.

El amplificador de potencia, por su parte, cumplió con los requisitos mandatorios de ancho de banda, potencia de salida y OIP3 establecidos en el anteproyecto, alcanzando un ancho de banda de 70~MHz, una potencia de salida de 30,92~dBm en las proximidades del P1dB y un OIP3 de 39,83~dBm. Los parámetros que no llegaron a satisfacer plenamente el objetivo de diseño, la frecuencia central y la PAE, resultaron sin embargo muy próximos a dicho objetivo: la frecuencia central medida de 2,11~GHz se ubicó a solo 90~MHz del valor de diseño de 2,2~GHz, mientras que la PAE de 30,51\,\% obtenida en las proximidades del P1dB se aproximó de manera considerable al 35\,\% especificado, incluso superándolo al operar el amplificador dentro de su región de compresión. Cabe destacar que dichos objetivos de diseño fueron originalmente formulados pensando en una arquitectura Doherty completa, compuesta por un amplificador de portadora y un amplificador de pico operando en conjunto, y no en un único amplificador Clase AB operando de manera aislada como el desarrollado a partir de la redefinición de alcance descripta en el Capítulo~\ref{sec:alcance}. Que un amplificador Clase AB simple haya logrado aproximarse en tal medida a especificaciones concebidas para una topología de mayor complejidad y de mayor eficiencia constituye un resultado alentador de cara a una futura implementación Doherty completa, en la cual cabe esperar que la combinación de ambas ramas de amplificación permita satisfacer los objetivos no alcanzados y superar holgadamente los que ya fueron cumplidos.

De manera análoga a lo observado en el acoplador, la comparación entre los resultados medidos y los simulados en ADS a nivel esquemático evidenció el origen de las principales discrepancias del amplificador, atribuidas a la inductancia parásita de los componentes discretos de la red de estabilidad y a una resonancia fuera de banda no capturada por el modelo esquemático original. La identificación puntual de estas causas a través de un proceso de análisis similar al del acoplador deja establecidas dos mejoras concretas y directamente accionables para una futura iteración del diseño: la incorporación de la inductancia equivalente serie de los componentes de la red de estabilidad desde las primeras etapas de simulación y un modelado más detallado de los efectos parásitos del layout físico y del empaquetado de los componentes.

El trabajo futuro cuenta con un conjunto de mejoras ya identificadas y fundamentadas que permitirían acercar aún más el desempeño de ambos circuitos a sus respectivos comportamientos ideales. Para el acoplador, como fue mencionado en la Sección~\ref{sec:resultados_acoplador}, un cambio de sustrato dieléctrico permitiría reducir aún más sus dimensiones físicas, mientras que la incorporación de un modelo tridimensional del conector SMA en la simulación electromagnética permitiría reducir la discrepancia entre los resultados simulados y los medidos. La asimetría observada entre los parámetros de reflexión de los puertos de salida del acoplador, sin una explicación igualmente clara hasta el momento, constituye asimismo un aspecto que amerita un estudio más profundo. Para el amplificador, además de las mejoras de modelado ya mencionadas, el análisis de radioenlace desarrollado en la Sección~\ref{sec:radioenlace} estableció que un incremento del margen disponible, ya sea mediante una mayor ganancia de antena o mediante una reducción de la temperatura de ruido del sistema receptor, habilitaría el empleo de esquemas de modulación de mayor orden, mejora que además beneficiaría de manera directa a los esquemas QPSK y 8-PSK ya viables, ampliando el margen operativo y la fracción de cada pasada satelital disponible para la comunicación.

Naturalmente, la continuación más significativa de este trabajo reside en completar la arquitectura Doherty originalmente planteada, incorporando el amplificador de pico y la red de combinación de salida correspondiente. El acoplador desarrollado y caracterizado en este trabajo se encuentra ya en condiciones de cumplir su función como red de división de señal de entrada en dicha arquitectura, mientras que el amplificador Clase AB caracterizado puede ser directamente reutilizado como amplificador de portadora del sistema Doherty completo. En este sentido, ambos desarrollos presentados no constituyen resultados aislados, sino los dos primeros bloques constitutivos de un sistema de mayor complejidad, cuya continuación resulta ampliamente respaldada por el desempeño individual alcanzado por cada uno.

Finalmente, más allá de los resultados técnicos particulares de cada circuito, este trabajo permitió atravesar de manera integral las etapas propias del desarrollo de un módulo de radiofrecuencia para aplicaciones espaciales, desde el diseño teórico y la simulación asistida por computadora hasta la fabricación de circuitos impresos y la caracterización experimental en laboratorio, incluyendo la articulación con instituciones externas para acceder a instrumental de medición especializado no disponible en la institución de origen. La necesidad de redefinir el alcance original del proyecto debido a restricciones de tiempo, costo y disponibilidad de equipamiento constituyó un aprendizaje relevante sobre la gestión de un proyecto de ingeniería real, en el cual las condiciones planteadas rara vez coinciden exactamente con la planificación inicial. La plataforma de diseño, simulación y caracterización desarrollada y validada a lo largo de este trabajo, junto con el conjunto de mejoras concretas identificadas para su continuación, sienta una base sólida y reutilizable para el desarrollo de futuros módulos de radiofrecuencia orientados a aplicaciones aeroespaciales.